% ---
% Capitulo de revisão de literatura
% ---
\begin{comment}
Definição de problemas multiobjetivo
O que é;
O que são métodos escalares e porque são usados
Funções de Benchmark
Funções --- 

\end{comment}

\chapter{Referencial Teórico}\label{cap:referencial_teorico}

Neste capítulo, serão apresentados conceitos fundamentais que foram essenciais para o desenvolvimento do software proposto neste documento. 
A compreensão desses conceitos foi crucial para embasar e contextualizar a elaboração e funcionalidades do software em questão.

\section{Transtorno do espectro autista (TEA)}

O TEA é um distúrbio que afeta a comunicação social, 
dificultando o aprendizado em certos assuntos, 
além de resultar em comportamentos repetitivos e restritivos \cite{bordini2024exploring}. 
Geralmente, os sintomas do autismo são perceptíveis antes dos três anos de idade. 
Além das dificuldades na comunicação, comportamento e interações sociais, 
em alguns casos, indivíduos com autismo podem apresentar padrões alimentares atípicos \cite{london2007role}, 
bem como habilidades excepcionais de memória e compreensão em determinadas áreas \cite{treffert2009savant}.

É crucial que o tratamento de crianças com TEA seja iniciado o mais precocemente possível \cite{bosa2006autismo}. 
Dentre as abordagens utilizadas no tratamento, 
destacam-se a ABA (Análise do Comportamento Aplicada) e o TEACCH 
(\textit{Treatment and Education of Autistic and related Communication-handicapped Children}). 
A abordagem TEACCH é direcionada às necessidades individuais da criança com TEA, 
considerando seus pontos fortes e interesses específicos \cite{seiford2007perceived}.

\section{Análise de Comportamento aplicados}\label{otimizacao}

O \textit{Applied Behavior Analysis} (ABA) é uma abordagem terapêutica que envolve um terapeuta e um paciente, 
sentados em lados opostos de uma mesa. Durante as sessões de ABA, 
o terapeuta apresenta uma condição ou estímulo (entrada), 
enquanto o paciente responde com um comportamento específico (saída) \cite{artoni2011accessible}.

Se o comportamento do paciente é considerado adequado e alinhado com os objetivos terapêuticos, 
o terapeuta recompensa o paciente com um reforço positivo. 
Isso pode incluir atividades como cantar uma canção, 
comemorar o sucesso ou fornecer um brinquedo que seja especialmente gratificante para a criança \cite{artoni2018technology}.

No entanto, se o comportamento do paciente não atende aos critérios estabelecidos, 
o terapeuta intervém, aplicando outras condições ou estratégias terapêuticas para auxiliar o paciente a desenvolver comportamentos mais apropriados e desejados \cite{artoni2011accessible}.

Além do ABA, existem outros métodos amplamente conhecidos, como o PECS 
(\textit{Picture Exchange Communication System}) e o TEACCH 
(\textit{Treatment and Education of Autistic and related Communication-handicapped Children}). 
O PECS é um método simples que utiliza imagens para estabelecer uma comunicação funcional, 
principalmente voltado para autistas não verbais. 
Já a estratégia TEACCH é aplicada com foco na adaptação, 
comunicação e organização da criança. Enquanto isso, 
o ABA concentra-se mais no comportamento e na aprendizagem geral da criança \cite{viana2020autismo}.

\section{Treino por Tentativas Discretas}

O \textit{Discrete Trial Training} (DTT), 
também conhecido como Treinamento por Tentativas Discretas, 
é uma técnica de ensino amplamente empregada no âmbito da Análise do Comportamento Aplicada (ABA), 
especialmente no tratamento de indivíduos com TEA. 
No DTT, uma série de provas, treinos e atividades são estruturados de forma a permitir que a criança aprenda por meio de um processo de tentativa e erro organizado e focado \cite{artoni2012designing}.

No DTT, o estímulo pode ser um cartão 
(\textit{card} regularmente feito a mão) ou um brinquedo, 
nos quais o terapeuta apresenta um comando específico. 
Quando a criança realiza o comportamento desejado em resposta ao comando, 
o terapeuta aplica um reforço, recompensando a criança por sua resposta positiva \cite{artoni2011accessible}.

\section{Programas do Treino por Tentativas Discretas}

Cada indivíduo com TEA requer abordagens terapêuticas personalizadas, 
e os terapeutas podem escolher entre diferentes programas terapêuticos, que incluem:

\begin{itemize}
    \item \textbf{\textit{Matching} (Correspondência - Imagem para Imagem):}
    Nesse programa, um conjunto de estímulos visuais é apresentado, 
    e ao paciente é fornecido um estímulo idêntico ou correspondente a um dos apresentados anteriormente. 
    Um comando é emitido, e espera-se que o paciente mova o estímulo fornecido para a posição correspondente na mesa\cite{artoni2013portable}.

    \item \textbf{\textit{Receptive} (Receptivo - Palavra para Imagem):} 
    Neste esquema, um conjunto de estímulos é apresentado, 
    e o terapeuta emite um comando, como por exemplo, 'Toque na maçã'. 
    O paciente é instruído a selecionar o estímulo correspondente à palavra indicada. 
    Esse programa tem por objetivo aprimorar a compreensão de comandos verbais e a habilidade de associar palavras a objetos ou imagens \cite{Antunes2012}.

    \item \textbf{\textit{Expressive} (Expressivo - Imagem para Palavra): } 
    Neste módulo, o terapeuta aponta para um estímulo e emite um comando, 
    como por exemplo, 'Diga o nome'. 
    O paciente é encorajado a verbalizar o nome do objeto representado no estímulo. 
    Esse programa tem por objetivo melhorar as habilidades expressivas e de comunicação verbal do paciente 
    \cite{artoni2011accessible}.
\end{itemize}

\section{Níveis de Dificuldade nos Treinos do DTT}

À medida que o tratamento avança, 
o indivíduo passa por uma progressão de níveis de dificuldade. 
Cada um desses níveis adota uma abordagem distinta e abarca uma variedade de cenários de erro, 
com o intuito de aprimorar continuamente as habilidades terapêuticas. 
Nesta seção, são apresentadas informações detalhadas sobre cada um desses níveis de dificuldade \cite{artoni2011accessible}.

\begin{itemize}
    \item \textbf{Treino em massa:} 
    Este nível concentra-se em um único estímulo para assegurar o sucesso da criança. 
    Caso a criança execute a ação desejada, o terapeuta pode intervir, 
    movendo a mão da criança para a ação necessária ou fornecendo uma sugestão.
    
    \item \textbf{Fase de distração:} 
    Aqui, são introduzidos estímulos adicionais para cenários de erro. 
    Inicialmente, são apresentados estímulos neutros, 
    seguidos por estímulos não neutros, 
    os quais podem ser mais chamativos ou distrativos.

    \item \textbf{Treino estendido:} 
    Neste estágio, os estímulos usados nos cenários de erro são aqueles previamente dominados pelo paciente, 
    promovendo um desafio adicional para consolidar as habilidades aprendidas.

    \item \textbf{Rotação aleatória:} 
    Em caso de erro, a intervenção não se baseia em dicas. 
    Nos cenários de acerto, o paciente é reforçado positivamente. 
    Em ambos os cenários, os estímulos são embaralhados e novos estímulos são introduzidos de maneira aleatória na mesa.
\end{itemize}

\section{Trabalhos Relacionados}

Na literatura, existem textos que propõem softwares para abordar desafios associados ao método ABA, 
como ilustrado em trabalhos de Artoni mencionados na tabela \ref{tab:trabalhos-relacionados}.

\begin{table}[h]
    \caption{Trabalhos Relacionados}
    \centering
    \begin{tabular}{|p{4cm}|p{3.40cm}|p{3.40cm}|p{3.40cm}|}
    \hline
    \textbf{Descrição} & \textbf{Vantagens} & \textbf{Desvantagens} & \textbf{Semelhanças} \\
    \hline
        Desenvolvimento de um \textit{software} para auxiliar no aprendizado de comunicação alternativa de crianças com autismo \cite{artoni2011accessible}. 
        &
        Um \textit{software} voltado principalmente para comunicação é mais especifico no que é proposto.
        &
        Requer treinamento ou familiarização com o sistema para garantir seu uso adequado.
        &
        É um \textit{software} planejado para uma área específica do ABA e todas as funcionalidades voltadas para ela.
        \\
    \hline
        Criação de uma interface \textit{web} baseada em DTT para o ensino de formas, cores e números em terapias \cite{artoni2012designing}. &
        Facilidade de acesso e autonomia da criança ao realizar os treinos. &
        Acessar e interagir com uma página \textit{web} envolve um processo mais complexo em comparação à interação direta com um aplicativo. &
        É um \textit{software} voltado para DTT onde visa o ensino da criança por meio de treinos. \\
    \hline
        Aplicativo \textit{mobile} para anotações rápidas, visando o método ABA \cite{artoni2013portable}. &
        Anotação de forma rápida e compartilhada com a família. &
        Sistema exclusivamente voltado para anotações, sem qualquer interação direta com os pacientes. &
        \textit{Software} voltado à coleta de dados durante as sessões de terapia. \\
    \hline
    \end{tabular}
    \label{tab:trabalhos-relacionados}
    {
        Fonte: Elaborada pelo autor
    }
\end{table}
